\section{Ideas}
\subsection{Общие:}
\par если задача на граф с возможностью ожидать в вершине, то можно добавить фиктивные ребра с весом 1, которые как бы будут симулировать ожидание (можно и не в саму себя).
\par если в дейкстре n ребер из вершин то не делать сет или пку а просто за линию выбрать минимум из не used (оптимайз).
\par если ТЛ, то заменять set на pq с alive/used.
\par если надо топK элементов, но на их порядок все равно то можно делать квик селект (nth_element) за линию, вместо сортировки за лог.
\par траенгуляция выпуклого многоугольника -- дерево (всегда $n - 2$ треугольника для $n$ вершин).
\par если есть массив $a$ и число $k$, а так же массив $b$, где $bi = ai \% k$, то $\text{summ}(a) - \text{summ}(b) = k \cdot ()$.
\par думать в инверсии в тасках на сортировку (посмотреть на то, какие операции как меняют позиции элементов (мб кол-во какого-то элемента на четных / нечетных позициях)).
\par бинапы на функциях.
\par искать ноп для массивов без повторок за $\mathcal{O}(n \log m)$ путем выпиcывания для 2го массива вместо элементов их индексы в первом массиве, после чего делать поиск нвп.
\par поиск самого частого элемента на отрезке при помощи МО (тот разбор которого я писал).
\par если хотим найти минимум/максимум по весам на ребрах по оптимальным путям от корня до всех вершин, то просто надо взять минимум/максимум по максимальному/минимальному ребру которое входит в v для каждого v.
\par в каждой позиции строки заканчиваетя не больше $\mathcal{O}(\sqrt{n})$ различных шаблонов т.к. сумммарная длинна у них.
\subsection{Геома:}
\par посчитать площадь объединения многоугольников -- выписать все $x$ (включая пересечения сторон разных многоугольников) и в промежутках между ними (это много трапеций, поэтому  $(\text{len}(l) + \text{len}(r)) / 2 * (xr - xl)$).
\par просто площадь выпуклого многоугольника -- сумма площадей треугольников из p0.
\par найти ближайшие точки -- разделяйка по x, сортим внутри половины по y и ищем новые в полоске вокруг центра (для каждой точки перебор костанты точек)
для поиска треугольника с мин макс стороной такой же алгоритм.
\subsection{ДП:}
\par если хочется считать дп с обновлениями в точке, которые будут вести изменения на отрезке, то можно делать ДО , в вершине хранить небольшой отрезок ДП и мержить.
\par бфс по таблице, если там дп с простыми переходами (попасть из (1, 1) в (n, m)).
\par пересчитывать треугольники перебирая маску диагонали.
\par делать переливайку до/личао/дерева в дфс по детям.


% \textbf{Бинпоиск по ответу:} если «подходит $x$» монотонно по $x$ -- искать границу; проверку писать жадно или паросочетанием. \textbf{Параллельный бинпоиск:} много запросов «первый момент, когда выполнено» -- бинпоиск каждого, но проверки для всех запросов делаются одним проходом по моментам.
% \par \textbf{Оффлайн:} если запросы не зависят от ответов -- отсортировать и обработать в удобном порядке (по правой границе + Фенвик; Mo; дерево отрезков по времени + откат СНМ для динамической связности). \textbf{Разделяй и властвуй:} посчитать вклад пар через середину, рекурсивно левую и правую половины; на деревьях -- центроиды.
% \par \textbf{Meet-in-the-middle:} $n\le40$ -- разбить на две половины по $2^{20}$, посчитать множества и склеить сортировкой или двумя указателями. \textbf{Битсет:} перебор вида «есть ли сумма $s$» ускоряется в 64 раза сдвигами \texttt{bitset}.
% \par \textbf{Корневая декомпозиция:} блок размера $\sqrt n$: обновление за $O(\sqrt n)$, запрос за $O(1)$ или наоборот; «тяжёлые/лёгкие» по степени: для тяжёлых ($>\sqrt m$) таблицы, для лёгких -- перебор.
% \par \textbf{Хеширование:} модуль $2^{61}-1$ со случайной базой (не $2^{64}$ -- ломается Туэ--Морса); для деревьев -- хеш поддерева как функция от отсортированных хешей детей. \textbf{Рандом:} равновероятный выбор вместо перебора, отжиг/hill climbing для оптимизационных задач без структуры.
% \par \textbf{Деление с округлением.} В C++ \texttt{a/b} округляет к нулю. $\lfloor a/b\rfloor$ при $a<0$: \texttt{a/b - (a\%b!=0 \&\& ((a\%b<0)!=(b<0)))}; $\lceil a/b\rceil=\lfloor(a+b-1)/b\rfloor$ при $a\ge0,b>0$. Остаток в $[0,m)$: \texttt{((a\%m)+m)\%m}.
